% ------------------------------------------------------------------------
% file `3_uaa1_fonctions_2022-2023-appliquer-5-solution-body.tex'
%   in folder `xsim/'
%
%     solution of type `appliquer' with id `5'
%
% generated by the `solappliquer' environment of the
%   `xsim' package v0.21 (2022/02/12)
% from source `3_uaa1_fonctions_2022-2023' on 2022/10/25 on line 185
% ------------------------------------------------------------------------
    1 point pour $m$ et 1 point pour $p$.\\
    -0,5 point si l'expression entière n'est pas notée.
    \begin{enumerate}
        \item
              \begin{tblr}[t]{columns={leftsep=0pt},
                      row{1}={font=\bfseries},
                      colspec={XX},
                      cell{3}{1}={r=1,c=2}{c}
                  }
                  Déterminons m                         & Déterminons p        \\
                  $m = \dfrac{12}{-8} = -\dfrac{3}{2}$  &
                  $(-2;8)\in f_1 \Rightarrow \begin{aligned}[t]
                                                     8 & = -\dfrac{3}{2}\cdot (-2) + p \\
                                                     8 & = 3 + p                       \\
                                                     p & = 11                          \\
                                                 \end{aligned}$ \\
                  \fbox{$f_1(x) = -\dfrac{3}{2}x + 11$} &
              \end{tblr}
        \item
              \begin{tblr}[t]{columns={leftsep=0pt},
                      row{1}={font=\bfseries},
                      colspec={XX},
                      cell{3}{1}={r=1,c=2}{c}
                  }
                  Déterminons m                      & Déterminons p \\
                  $m = \dfrac{14}{4} = \dfrac{7}{2}$ &
                  $f_2$ est linéaire $\Rightarrow p = 0$             \\
                  \fbox{$f_2(x) = \dfrac{7}{2}x$}    &
              \end{tblr}
        \item
              \begin{tblr}[t]{columns={leftsep=0pt},
                      row{1}={font=\bfseries},
                      colspec={XX},
                      cell{3}{1}={r=1,c=2}{c}
                  }
                  Déterminons m                              & Déterminons p \\
                  $G_{f_3} \parallel G_g \Rightarrow m = -5$ &
                  $(-2;0)\in f_3 \Rightarrow \begin{aligned}[t]
                                                     0 & = -5 \cdot (-2) + p \\
                                                     0 & = 10 + p            \\
                                                     p & = -10               \\
                                                 \end{aligned}$         \\
                  \fbox{$f_3(x) = -5x - 10$}                 &
              \end{tblr}
        \item
              Les deux points ont la même ordonnée $\Rightarrow$ la fonction est constante.\\
              \fbox{$f_4(x) = -3$}
        \item
              Les points $(0,1)$ et $(4,0)$ appartiennent à $f_5$ par lecture graphique.\par
              \begin{tblr}[t]{columns={leftsep=0pt},
                      row{1}={font=\bfseries},
                      colspec={XX},
                      cell{3}{1}={r=1,c=2}{c}
                  }
                  Déterminons m                    & Déterminons p                  \\
                  $m=-\dfrac{1}{4}$                & $(0,1)\in f_5 \Rightarrow p=1$ \\
                  \fbox{$f_5(x)=-\dfrac{1}{4}x+1$} &
              \end{tblr}
    \end{enumerate}
